\documentclass[a4paper,12pt]{report}

\usepackage[a4paper]{geometry}
\usepackage[utf8]{inputenc}
\usepackage{amssymb,amsmath,amsthm}
\usepackage{graphicx}
\usepackage{url}
\usepackage{epsfig}
\usepackage[italian]{babel}
\usepackage{setspace}
\usepackage{tesi}
\usepackage{import}
\usepackage[a-1b]{pdfx}
\usepackage{hyperref}
\usepackage{subcaption}

\begin{document}
% FRONTESPIZIO
\title{Analisi quantitativa e percettiva di video creati con generatori AI}
\author{Federico COSCIA}
\dept{Corso di Laurea in Informatica (L-31)} 
\anno{2023-2024}
\matricola{977772}
\relatore{Raffaella Lanzarotti}
\correlatore{Alessandro D'Amelio}

% CORPO
\import{Chapters/}{Preface.tex}
\import{Chapters/}{Introduzione.tex}
\import{Chapters/Capitolo1}{Capitolo1.tex}
\import{Chapters/Capitolo2}{Capitolo2.tex}
\import{Chapters/Capitolo3}{Capitolo3.tex}
\import{Chapters/Capitolo4}{Capitolo4.tex}
\import{Chapters/}{Conclusioni.tex}

\phantomsection		% da usare insieme a addcontentsline altrimenti i bookmark nel pdf puntano alle pagine sbagliate
\addcontentsline{toc}{chapter}{\bibname}
\bibliographystyle{plain}
\bibliography{bibliography.bib}
\import{Chapters/}{Ringraziamenti.tex}
\end{document}